\begin{table}[t]
\centering
\small
\begin{tabular}{lrrl}
\toprule
Correlation against forecast skill & $r$ & $n$ & leave-one-out range \\
\midrule
\emph{Panel A --- 18 BDG2 buildings, three countries} &  &  &  \\
\quad cooling-degree-day share & +0.010 & 18 & -0.308 to +0.077 \\
\quad controllable (HVAC) fraction & +0.283 & 18 & +0.029 to +0.378 \\
\quad mean outdoor temperature & -0.012 & 18 & -0.266 to +0.051 \\
\quad median load & -0.014 & 18 & -0.189 to +0.120 \\
\midrule
\emph{Panel B --- 6 Chinese provinces, independent replication} &  &  &  \\
\quad cooling-degree-day share & +0.054 & 6 & -0.377 to +0.582 \\
\quad mean outdoor temperature & +0.216 & 6 & -0.290 to +0.782 \\
\quad load--temperature correlation & -0.060 & 6 & -0.619 to +0.713 \\
\midrule
\emph{controllable fraction} vs climate (not vs skill) & \textbf{+0.489} & 18 & -- \\
\bottomrule
\end{tabular}
\caption{Forecast skill is uncorrelated with climate, controllable fraction, temperature and load size, and every leave-one-out range straddles or nearly straddles zero --- which at these sample sizes is what a null looks like. Panel B is the replication on a panel the null was not fitted to. The final row is the one relationship in the study that is real, and note what it is between: available flexibility tracks climate strongly, while predictability does not track anything. Skill and flexibility are independent deployment axes.}
\label{tab:correlations}
\end{table}
